% =====================================================================
\section{Identificación y Contextualización}
% =====================================================================

% ---------------------------------------------------------------------
\subsection{Contextualización de la empresa}
% ---------------------------------------------------------------------

Desarrollo de Tecnologías y Sistemas SpA (DTS) es una empresa chilena de electrónica con casa matriz en la comuna de Macul, Santiago. Se constituyó el 11 de enero de 1991 y se transformó en sociedad por acciones en noviembre de 2020. Su giro principal es la industria electrónica aplicada a los mercados de defensa, seguridad y aeroespacial, y sus principales clientes son las Fuerzas Armadas chilenas. La empresa registra ventas anuales entre 100.000 y 350.000 unidades de fomento (UF) y una dotación de entre 50 y 250 trabajadores \parencite{redconecta}, cifras que la ubican en el tramo de gran empresa según el criterio de ventas del Estatuto Pyme \parencite{minecon}.

Las áreas de negocio de DTS abarcan el desarrollo y la modernización de sistemas, el mantenimiento de equipos electrónicos y electromecánicos, las calibraciones, el soporte de sistemas, la guerra electrónica, el mando y control, el control automático, los bancos de prueba, la fibra óptica militar y el entrenamiento y simulación. Además de su casa matriz, la empresa mantiene oficinas técnicas en Iquique, Antofagasta, Talcahuano y Punta Arenas \parencite{dts}.

El practicante trabajó en el Departamento de Investigación, Desarrollo e Innovación (I+D+i), un equipo de seis personas orientado al desarrollo de tecnologías militares. Dentro de ese departamento, el desarrollo de capacidades espaciales corresponde a una línea de trabajo incipiente, que el Plan Estratégico 2025--2029 contempla pero que hasta ese momento no tenía ningún proyecto formal que sirviera de referencia. El proyecto BIFROST es el primer desarrollo espacial de la empresa.


% ---------------------------------------------------------------------
\subsection{Descripción del proyecto y rol del practicante}
% ---------------------------------------------------------------------

El proyecto que desarrolló el practicante recibió el nombre BIFROST, elegido por él mismo como analogía del puente que comunica dos puntos lejanos. Su objetivo general es diseñar y validar, a nivel de arquitectura, una carga útil de comunicaciones del tipo \textit{Store \& Forward} para un satélite CubeSat 12U, capaz de recibir mensajes autónomos desde estaciones remotas terrestres en banda UHF, almacenarlos temporalmente a bordo y descargarlos con posterioridad a estaciones terrenas ubicadas en territorio nacional.

La carga útil funciona como un repetidor no instantáneo. Un receptor en la banda de 401 a 403 MHz escucha las transmisiones que las plataformas de recolección de datos (DCP) emiten de forma autónoma, sin coordinación previa con el satélite. El sistema compensa el corrimiento en frecuencia que introduce el efecto Doppler, demodula y valida cada mensaje, y guarda los mensajes correctos en memoria no volátil junto con su marca temporal. Cuando el satélite alcanza una ventana de visibilidad sobre una estación terrena, y bajo comando del computador de a bordo, la carga útil descarga la información almacenada mediante un enlace en banda UHF. Además, la carga útil intercambia telemetría, telecomandos y estado operativo con el resto del satélite a través del bus CAN de la plataforma.

La práctica se desarrolló entre el 2 de enero y el 28 de febrero de 2026, en jornada completa, y cubrió las tres primeras etapas del ciclo de vida del proyecto \parencite{nasaseh}. La Figura \ref{img:gantt} presenta la carta Gantt que el practicante elaboró durante ese período. El área sombreada corresponde al trabajo que él realizó y el resto, a la planificación que el equipo heredó para 2026 y 2027, mientras que los rombos marcan los hitos de revisión formal. Durante la práctica el proyecto cerró dos de esos hitos, que fueron la Revisión de Requerimientos del Sistema (SRR), el 30 de enero de 2026, y la Revisión de Diseño Preliminar (PDR), el 27 de febrero de 2026. La Tabla \ref{tab:hitos} resume ambos hitos junto con los entregables que el practicante generó en cada uno de ellos.

El practicante se desempeñó como Ingeniero Civil Eléctrico en práctica profesional y estuvo a cargo de elaborar la documentación técnica de las dos revisiones. Su supervisor directo fue Juan Carlos Sepúlveda, Ingeniero de Desarrollo Espacial, y la supervisión técnica estuvo a cargo de Álvaro Valenzuela, doctor en sistemas espaciales. Ambos reportan a Lino Acuña, jefe del Departamento I+D+i, quien aprobó los entregables. La única contraparte externa del proyecto fue INVAP, empresa argentina que participó como asesor técnico en el marco del acuerdo de cooperación ENAER--DTS--INVAP suscrito en marzo de 2025, y que revisó el avance del diseño en cuatro instancias formales a lo largo de la práctica.

\fig[gantt]{Carta Gantt del proyecto BIFROST. Elaboración propia.}{width=0.8\textwidth}{gantt.png}

% [Tabla 1.1] Resume plazos, hitos y entregables. Cubre de una sola vez los
% criterios "plazos" y "alcance" de la rúbrica. Se menciona en el texto.
\begin{table}[H]
\centering
\begin{threeparttable}
\caption{Hitos de revisión y entregables generados durante la práctica.
Elaboración propia.\vspace{1ex}}
\label{tab:hitos}
\small
\begin{tabularx}{\textwidth}{|@{\hspace{4pt}}l l X@{\hspace{4pt}}|}
\toprule
\textbf{Fase} & \textbf{Hito y fecha} & \textbf{Entregables generados} \\
\midrule
Pre-Phase A & \textit{Kick-off} & Concepto de Operación (CONOPS) preliminar y
diagrama de bloques inicial del sistema \\

Phase A & SRR & CONOPS consolidado, 13 requerimientos de usuario, 21 requerimientos de sistema, matriz de trazabilidad, matriz de riesgos y carta Gantt \\

Phase B & PDR & Modelo MBSE en Capella con las cuatro capas ARCADIA, herramienta de simulación de enlace, matriz de ensayos y estimación ROM de costos \\
\bottomrule
\end{tabularx}
\begin{tablenotes}[flushleft]\footnotesize
\item Nota. El \textit{kick-off} se realizó el 16 de enero de 2026, la SRR el 30 de enero de 2026 y la PDR el 27 de febrero de 2026.
\end{tablenotes}
\end{threeparttable}
\end{table}


% ---------------------------------------------------------------------
\subsection{Objetivo de la empresa y problema a resolver}
% ---------------------------------------------------------------------

El proyecto BIFROST es el primer esfuerzo de DTS por desarrollar capacidades espaciales y responde a dos objetivos estratégicos de la empresa. El primero fue dar el paso formal inicial en esa línea de trabajo, que el Plan Estratégico 2025--2029 contempla pero que hasta ese momento no tenía ejecución. El segundo fue aprovechar la ventana de transferencia tecnológica que abrió el acuerdo con INVAP, capitalizando la experiencia de la contraparte en sistemas satelitales de recolección de datos. Ambos objetivos convergen en un mismo resultado esperado, que consiste en llevar el diseño de la carga útil hasta la madurez que exige la Revisión de Diseño Preliminar, de manera de demostrar su viabilidad técnica y arquitectónica antes de comprometer recursos en el diseño detallado y la fabricación.

El principal obstáculo para avanzar hacia esos objetivos era la capacidad disponible en el departamento. El equipo de I+D+i mantenía su dotación comprometida en proyectos de tecnología militar en curso, por lo que la empresa no contaba con personal libre para levantar una línea de trabajo nueva sin desatender la carga existente. Por esta razón, DTS decidió iniciar la línea espacial con el equipo mínimo que podía destinar, es decir, un ingeniero de planta con experiencia previa en desarrollo espacial y un practicante disponible durante los dos meses que dura la práctica profesional.

Para abordar esa restricción, la empresa incorporó al practicante como recurso dedicado al proyecto, bajo la conducción técnica del supervisor. El supervisor definió el marco de trabajo, esto es, un modelo de desarrollo secuencial y los hitos que el proyecto debía cumplir. El practicante asumió la ejecución, para lo cual investigó en detalle los entregables y los criterios de salida que corresponden a cada revisión \parencite{ieee15288-2}, seleccionó y operó las herramientas de modelado y simulación, y elaboró la totalidad de la documentación técnica del proyecto siguiendo el método ARCADIA \parencite{roques}.

De esta manera, el problema que aborda el presente informe se plantea desde la empresa y se puede resumir en la siguiente pregunta: ¿cómo puede DTS iniciar el desarrollo de su línea espacial, y llevarla hasta un nivel de madurez que justifique su continuidad, contando solamente con un ingeniero de planta y un practicante disponible durante dos meses? Las secciones siguientes desarrollan la metodología que el proyecto adoptó y los resultados que obtuvo.


% ---------------------------------------------------------------------
\subsection{Antecedentes, causas e importancia de su resolución}
% ---------------------------------------------------------------------

El problema planteado surge de dos antecedentes. El primero es la baja madurez del desarrollo espacial, tanto en DTS como en la industria nacional. La empresa acumula más de tres décadas de experiencia en electrónica, aviónica y mantenimiento de sistemas, pero no había aplicado ninguna de esas capacidades a un sistema destinado a operar en órbita. Como consecuencia, el practicante tuvo que generar la documentación y las herramientas de análisis sin precedentes internos que le sirvieran de referencia, y tuvo que justificar cada criterio de diseño a partir de la norma o de la literatura técnica, ya que la empresa no contaba con experiencia acumulada en proyectos anteriores de este tipo.

El segundo antecedente es el acuerdo de cooperación que ENAER, DTS e INVAP suscribieron en marzo de 2025, y que habilitó formalmente la transferencia de conocimientos y tecnología en materia espacial y de radares. Ese acuerdo aportó una contraparte con experiencia acreditada en sistemas de recolección de datos satelitales, que revisó el diseño en instancias formales y fijó el contexto dentro del cual DTS debía mostrar avances concretos.

La resolución de este problema es importante por tres razones. En primer lugar, para DTS, porque el resultado determina si la línea espacial continúa y porque deja a la empresa una metodología y unas bases técnicas que antes no tenía. En segundo lugar, para el país, ya que el proyecto aporta al ecosistema espacial que impulsan el Centro Espacial Nacional y la Fuerza Aérea de Chile. Por último, para el usuario final, porque la geografía chilena deja zonas extensas sin cobertura terrestre confiable, donde un canal satelital de respaldo tiene un valor operacional directo.
